As we saw in \S~28, the values of $l$ and $m$ determine the angular
dependence of a particle's wave function and, with it, all the rotational
symmetry properties of that function. In their most general form these
properties are specified by stating how wave functions transform when the
coordinate system is rotated.

The wave function $\psi_{LM}$ of a system of particles, with prescribed
angular momentum $L$ and projection $M$, is unchanged\footnote{Apart from an
inessential phase factor.} only by rotations of the coordinate system about
the $z$ axis. After any rotation that alters the direction of this axis, the
projection of the angular momentum on $z$ no longer has a definite value. In
the new coordinate axes the wave function must therefore become, in general,
a superposition (linear combination) of the $2L+1$ functions associated with
the possible values of $M$ at fixed $L$. Equivalently, rotations of the
coordinate system transform the $2L+1$ functions $\psi_{LM}$ into one
another\footnote{In mathematical terminology, these functions realize an
\emph{irreducible representation} of the rotation group. The number of
functions transformed into one another is called the \emph{dimension of the
representation}; it is understood that no choice of other linear
combinations of these functions can reduce this number.}. The transformation
law---that is, the superposition coefficients as functions of the angles
through which the coordi\SourcePageJoin{247}{250}nate axes are rotated---is
fixed completely by the value of $L$. Angular momentum thus takes on the
meaning of a quantum number that classifies the states of a system by their
transformation properties under rotations of the coordinate system. This
aspect of angular momentum is especially important in quantum mechanics
because it does not depend directly on an explicit angular dependence of the
wave functions: their mutual transformation law can be formulated by itself,
without referring to such a dependence.
